\subsection{Locked Internal Holdout}
After the six matched early and full reliability-aware checkpoints were fixed, the unchanged 837-image guard manifest was evaluated once without retraining, threshold tuning, or checkpoint reselection. The split inventory contains 2,287 boxes, while the standardized evaluator retains 1,264 ground-truth boxes after its label and geometry contract; both counts are reported to avoid conflating inventory and evaluator populations.

\tabref{tab:guard} shows that reliability-aware fusion improves mean AP50 from $0.9588\pm0.0058$ to $0.9674\pm0.0033$. The paired AP50 gain is $0.0086\pm0.0062$ and is positive for all three seeds. Mean F1 increases by $0.0069\pm0.0089$. AP75 improves slightly on average but is mixed by seed, including a negative seed-2 difference. The guard therefore provides useful locked internal holdout evidence, not an independent public test or a basis for post-hoc tuning.

\begin{table*}[t]
\centering
\caption{Locked 837-image internal holdout. Project-local AP50/AP75 and F1 are reported after all six checkpoints were fixed.}
\label{tab:guard}
\small
\begin{tabular}{@{}lrrr@{}}
\toprule
Model/statistic & F1 & AP50 & AP75\\
\midrule
Matched early, mean $\pm$ SD & $0.9253\pm0.0170$ & $0.9588\pm0.0058$ & $0.9007\pm0.0209$\\
Reliability $p=0.15$, mean $\pm$ SD & $0.9322\pm0.0106$ & $0.9674\pm0.0033$ & $0.9029\pm0.0323$\\
Paired difference, mean $\pm$ SD & $+0.0069\pm0.0089$ & $+0.0086\pm0.0062$ & $+0.0022\pm0.0173$\\
Positive seed pairs & 2 / 3 & 3 / 3 & 2 / 3\\
\bottomrule
\end{tabular}
\end{table*}

\subsection{Efficiency}
The reliability-aware system adds 1,684 parameters and approximately 0.77 GFLOPs under the reported profiling scope. Measurement uses batch size 1, image size 640, model evaluation mode, \texttt{torch.inference\_mode()}, 200 warm-up iterations, five trials of 1,000 timed iterations, and CUDA synchronization. Raw-forward timing covers the backbone call only; detector inference covers the torchvision FCOS transform, backbone, head, NMS, and postprocessing but excludes data loading and file I/O. \tabref{tab:efficiency} reports the mean of the median latency across the two fixed checkpoints and the observed peak-memory range. Reliability-aware fusion is slightly slower and uses more memory; it should not be described as faster.

\begin{table*}[t]
\centering
\caption{Efficiency on RTX 3090, batch size 1, image size 640. Latency is the mean of the median latency from two checkpoints. Raw-forward and detector-inference boundaries are reported separately.}
\label{tab:efficiency}
\small
\begin{tabular}{@{}lrrr@{}}
\toprule
Measurement & Matched early & Reliability $p=0.15$ & Difference\\
\midrule
Parameters & 6,591,609 & 6,593,293 & +1,684\\
Raw-forward GFLOPs & 24.44 & 25.22 & +0.77\\
Detector-inference GFLOPs & 105.21 & 105.98 & +0.77\\
Raw-forward latency (ms) & 10.43 & 10.61 & +0.18\\
Detector-inference latency (ms) & 22.89 & 23.30 & +0.41\\
Raw-forward peak CUDA memory (MB) & 199.7--222.4 & 354.6--374.9 & higher\\
Detector-inference peak CUDA memory (MB) & 207.5--229.6 & 361.7--382.7 & higher\\
\bottomrule
\end{tabular}
\end{table*}

\subsection{Cross-Dataset Evaluation on MM-UAV}
\label{sec:mmuav}
MM-UAV \cite{xu2025mmuav} is used to study transfer across a different RGB-infrared-event acquisition system. The authorized local research subset contains a frozen 7,187-row training manifest and a 1,845-row devval manifest spanning 85 sequences; the devval annotations contain 4,198 vehicle boxes. Because these sequences were exposed during prior engineering, the split is not treated as an independent blind test. All MM-UAV metrics use COCO AP@[0.50:0.95], AP50, AP75, AR1, AR10, and AR100, which are distinct from the project-local TriAir AP50 and AP75 metrics reported above. Dataset media and annotations are not redistributed, and this manuscript does not assert public redistribution rights.

\subsubsection{Frozen zero-shot normalized-grid stress test}
Six frozen TriAir checkpoints---matched early fusion and reliability-aware $p=0.15$ for seeds 0, 1, and 2---are evaluated exactly once on all 1,845 devval rows. RGB, infrared, and event observations are decoded independently, letterboxed to $640\times640$, and concatenated as three RGB channels plus one infrared and one event channel. Ground-truth boxes follow RGB geometry. The adapter is deterministic and parameter free, but it does not establish physical pixel correspondence across sensors. The evaluator uses a score threshold of 0.001, NMS threshold of 0.6, and at most 100 detections per image.

All six checkpoints produce finite predictions for every image, with 184,500 valid decoded predictions per checkpoint, yet AP is effectively zero (\tabref{tab:mmuav-zero-shot}). This rules out an empty-output or numerical-failure explanation and instead indicates that the predicted geometry does not match the target annotations under the naive unregistered adapter. The reliability-minus-early differences are also effectively zero in absolute terms and are not interpreted as a robustness gain.

\begin{table*}[t]
\centering
\caption{Frozen TriAir checkpoints on exposed MM-UAV devval using the naive normalized-grid adapter. Metrics are COCO-style. The adapter does not establish physical cross-modal registration.}
\label{tab:mmuav-zero-shot}
\scriptsize
\begin{tabular}{@{}llrrrrrr@{}}
\toprule
Model group & Seed & AP & AP50 & AP75 & AR1 & AR10 & AR100\\
\midrule
Matched early fusion & 0 & $1.71{\times}10^{-8}$ & $8.54{\times}10^{-8}$ & 0 & 0 & 0 & $4.76{\times}10^{-5}$\\
Matched early fusion & 1 & $1.14{\times}10^{-8}$ & $1.14{\times}10^{-7}$ & 0 & 0 & 0 & $2.38{\times}10^{-5}$\\
Matched early fusion & 2 & 0 & 0 & 0 & 0 & 0 & 0\\
Reliability-aware $p=0.15$ & 0 & $8.67{\times}10^{-8}$ & $5.15{\times}10^{-7}$ & 0 & 0 & 0 & $1.91{\times}10^{-4}$\\
Reliability-aware $p=0.15$ & 1 & $2.08{\times}10^{-6}$ & $1.04{\times}10^{-5}$ & 0 & $4.76{\times}10^{-5}$ & $4.76{\times}10^{-5}$ & $1.19{\times}10^{-4}$\\
Reliability-aware $p=0.15$ & 2 & 0 & 0 & 0 & 0 & 0 & 0\\
\bottomrule
\end{tabular}
\end{table*}

\subsubsection{Supervised alignment-aware transfer benchmark}
The supervised benchmark tests whether target-domain training and explicit feature alignment recover detection performance, and whether TriAir initialization adds value beyond the matched target-only control. All variants use independent RGB, infrared, and event stems, deterministic modality decoding, learned feature-level alignment, Softplus box-distance output, a 128-channel FPN, and the same detector head. Raw five-channel concatenation is not used.

Exactly nine runs are trained: \textit{Scratch Equal}, \textit{TriAir Init Equal}, and \textit{TriAir Init Reliability}, each with seeds 0, 1, and 2. Every run uses $640\times640$ inputs, batch size 1, 10 epochs, 71,870 optimizer steps, AdamW with learning rate $10^{-4}$ and weight decay $10^{-4}$, a 500-step warmup, and cosine decay to $10^{-6}$. Augmentation, early stopping, devval monitoring during training, threshold tuning, and checkpoint selection are disabled. Only the final checkpoint is evaluated once on all 1,845 devval rows. For the initialized variants, only exact name-and-shape-compatible shared RepViT, FPN, and FCOS tensors are copied. The reported 99.20\% is the fraction of destination parameter and buffer elements copied by this exact rule; it is not a sample match rate or an architecture-identity claim. MM-UAV-specific input projection, stems, alignment modules, reliability scorer, and other unmatched tensors retain the matched seed-specific initialization.

\tabref{tab:mmuav-transfer} reports the independently verified three-seed summary. Supervised alignment-aware training raises AP from 0.000 under frozen naive-grid transfer to $0.2210\pm0.0030$ for Scratch Equal. TriAir Init Equal improves AP to $0.2340\pm0.0020$, and TriAir Init Reliability improves it further to $0.2503\pm0.0025$ while also attaining the highest AP50, AP75, and AR100 means. The complete seed-level records are reported in \tabref{tab:mmuav-transfer-seeds}.

\begin{table*}[t]
\centering
\caption{Corrected three-seed MM-UAV transfer summary on the 0--1 scale. AP is mean $\pm$ sample standard deviation; the remaining columns are three-seed means. Comparisons are descriptive and supervised. The best supervised row is bold.}
\label{tab:mmuav-transfer}
\scriptsize
\begin{tabularx}{\textwidth}{@{}YrrrrY@{}}
\toprule
Training setting & AP & AP50 & AP75 & AR100 & Interpretation\\
\midrule
Frozen TriAir, naive-grid zero-shot & 0.000 & 0.000 & 0.000 & 0.000 & Direct transfer failed\\
MM-UAV Scratch Equal & $0.2210\pm0.0030$ & 0.5567 & 0.1347 & 0.3530 & Aligned supervised training recovered performance\\
TriAir Init Equal & $0.2340\pm0.0020$ & 0.5797 & 0.1510 & 0.3717 & Source-domain pretraining was beneficial\\
\textbf{TriAir Init Reliability} & $\mathbf{0.2503\pm0.0025}$ & \textbf{0.6077} & \textbf{0.1750} & \textbf{0.3920} & Reliability-aware fusion improved performance further\\
\bottomrule
\end{tabularx}
\end{table*}

\begin{figure}[t]
\centering
\includegraphics[width=\columnwidth]{figures/fig5_mmuav_transfer.pdf}
\caption{Corrected three-seed COCO AP on the supervised MM-UAV transfer benchmark. Error bars are sample standard deviations. TriAir initialization improves over scratch equal fusion, and the reliability-aware initialized variant has the highest mean AP.}
\label{fig:mmuav-transfer}
\end{figure}

